\setcounter{section}{1}
\setcounter{theorem}{5}
\setcounter{equation}{4}
\setcounter{figure}{2}
\section{Preliminaries}

\label{sec:prelim}

In this section we introduce some definitions and known facts about
polynomials and complexity classes.

%
\subsection{Approximation theory}
\label{s:approxprelim}
We will use several results from approximation theory,\ each of which has
previously been used (in some form) in other applications of the polynomial
method to quantum lower bounds. \ We start with the basic inequality of A.A.
Markov \cite{markov1890question}.

\begin{lemma}[Markov]
\label{markovlem}Let $p$\ be a real polynomial, and suppose that
\begin{equation}
\max_{x,y\in\left[ a,b\right] }\left\vert p\left( x\right) -p\left( y\right)
\right\vert \leq H.
\end{equation}
Then for all $x\in\left[ a,b\right] $, we have
\begin{equation}
\left\vert p^{\prime}\left( x\right) \right\vert \leq\frac{H}{b-a}\deg\left(
p\right) ^{2},
\end{equation}
where $p^{\prime}(x)$ is the derivative of $p$ at $x$.
\end{lemma}

We'll also need a bound that was explicitly stated by Paturi \cite{paturi},
and which amounts to the fact that, among all degree-$d$
polynomials that are bounded within a given range, the Chebyshev polynomials
have the fastest growth outside that range.

\begin{lemma}[Paturi]
\label{paturilem}Let $p$\ be a real polynomial, and suppose that $\left\vert
p\left( x\right) \right\vert \leq1$\ for all $\left\vert x\right\vert \leq1$
. \ Then for all $x\leq1+\mu$, we have
\begin{equation}
\left\vert p\left( x\right) \right\vert \leq\exp\left( 2\deg\left( p\right)
\sqrt{2\mu+\mu^{2}}\right) .
\end{equation}
\end{lemma}

We now state a useful corollary of \Cref{paturilem}, which says (in effect)
that slightly shrinking the domain of a low-degree real polynomial can only
modestly shrink its range.

\begin{corollary}
\label{paturicor}Let $p$\ be a real polynomial of degree $d$, and suppose
that
\begin{equation}
\max_{x,y\in\left[ a,b\right] }\left\vert p\left( x\right) -p\left( y\right)
\right\vert \geq H.
\end{equation}
Let $\varepsilon\leq\frac{1}{100d^{2}}$\ and $a^{\prime}:=a+\varepsilon
\left( b-a\right) $. \ Then
\begin{equation}
\max_{x,y\in\left[ a^{\prime},b\right] }\left\vert p\left( x\right) -p\left(
y\right) \right\vert \geq\frac{H}{2}.
\end{equation}
\end{corollary}

\begin{proof}
Suppose by contradiction that
\begin{equation}
\left\vert p\left(  x\right)  -p\left(  y\right)  \right\vert <\frac{H}{2}
\end{equation}
for all $x,y\in\left[  a^{\prime},b\right]  $. \ By affine shifts, we can
assume without loss of generality that $\left\vert p\left(  x\right)
\right\vert <\frac{H}{4}$\ for all $x\in\left[  a^{\prime},b\right]  $. \ Then
by \Cref{paturilem}, for all $x\in\left[  a,b\right]  $\ we have
\begin{equation}
\left\vert p\left(  x\right)  \right\vert <\frac{H}{4}\cdot\exp\left(
2d\sqrt{2\left(  \frac{1}{1-\varepsilon}-1\right)  +\left(  \frac
{1}{1-\varepsilon}-1\right)  ^{2}}\right)  \leq\frac{H}{2}.
\end{equation}
But this violates the hypothesis.
\end{proof}

We will also need a bound that relates the range of a low-degree polynomial
on a discrete set of points to its range on a continuous interval. \ The
following lemma generalizes a result due to Ehlich and Zeller \cite{ez} and
Rivlin and Cheney \cite{rc}, who were interested only in the case where the
discrete points are evenly spaced.

\begin{lemma}
\label{ezrclem}Let $p$\ be a real polynomial of degree at most $\sqrt{k}$,
and let $0=z_{1}<\cdots<z_{M}=k$\ be a list of points such that $
z_{i+1}-z_{i}\leq1$\ for all $i$ (the simplest example being the integers $
0,\ldots,k$). \ Suppose that
\begin{equation}
\max_{x,y\in\left[ 0,k\right] }\left\vert p\left( x\right) -p\left( y\right)
\right\vert \geq H.
\end{equation}
Then
\begin{equation}
\max_{i,j}\left\vert p\left( z_{i}\right) -p\left( z_{j}\right) \right\vert
\geq\frac{H}{2}.
\end{equation}
\end{lemma}

\begin{proof}
Suppose by contradiction that
\begin{equation}
\left\vert p\left(  z_{i}\right)  -p\left(  z_{j}\right)  \right\vert
<\frac{H}{2}
\end{equation}
for all $i,j$. \ By affine shifts, we can assume without loss of generality
that $\left\vert p\left(  z_{i}\right)  \right\vert <\frac{H}{4}$\ for all
$i$. \ Let
\begin{equation}
c:=\max_{x\in\left[  0,k\right]  }\frac{\left\vert p\left(  x\right)
\right\vert }{H/4}.
\end{equation}
If $c\leq1$, then the hypothesis clearly fails, so assume $c>1$. \ Suppose
that the maximum, $\left\vert p\left(  x\right)  \right\vert =\frac{cH}{4}$,
is achieved between $z_{i}$\ and $z_{i+1}$. \ Then by basic calculus, there
exists an $x^{\ast}\in\left[  z_{i},z_{i+1}\right]  $\ such that
\begin{equation}
\left\vert p^{\prime}\left(  x^{\ast}\right)  \right\vert >\frac{2\left(
c-1\right)  }{z_{i+1}-z_{i}}\cdot\frac{H}{4}\geq\frac{\left(  c-1\right)
H}{2}.
\end{equation}
So by \Cref{markovlem},
\begin{equation}
\frac{\left(  c-1\right)  H}{2}<\frac{cH/4}{k}\deg\left(  p\right)  ^{2}.
\end{equation}
Solving for $c$, we find
\begin{equation}
c<\frac{2k}{2k-\deg\left(  p\right)  ^{2}}\leq2.
\end{equation}
But if $c<2$, then $\max_{x\in\left[  0,k\right]  }\left\vert p\left(
x\right)  \right\vert <\frac{H}{2}$, which violates the hypothesis.
\end{proof}

We also use a related inequality due to Coppersmith and Rivlin~\cite
{coppersmith-rivlin} that bounds a polynomial on a continuous interval in
terms of a bound on a discrete set of points, but now with the weaker
assumption that the degree is at most $k$, rather than $\sqrt{k}$. \ This
gives a substantially weaker bound.

\begin{lemma}[Coppersmith and Rivlin]
\label{lem:coppersmith_rivlin} Let $p$ be a real polynomial of degree at
most $k$, and suppose that $\left\vert p(x)\right\vert \leq 1$ for all
integers $x\in \{0,1,\ldots ,k\}$. \ Then there exist universal constants $
a,b$ such that for all $x\in \lbrack 0,k]$, we have
\begin{equation}
\left\vert p(x)\right\vert \leq a\cdot \exp \left( b\,\deg (p)^{2}/k\right) .
\end{equation}
\end{lemma}


%
\subsection{Symmetric polynomials}

\para{Univariate symmetrizations.}
Our starting point is the well-known \textit{symmetrization lemma} of Minsky
and Papert \cite{mp} (see also Beals et al.\ \cite{bbcmw}\ for its
application to quantum query complexity), by which we can often reduce
questions about multivariate polynomials to questions about univariate ones.

\begin{lemma}[Minsky--Papert symmetrization]
\label{symlem} Let $p:\left\{ 0,1\right\} ^{N}\rightarrow\mathbb{R}$\ be a
real multilinear polynomial of degree $d$, and let $q:\{0,1,\ldots,N\}\to
\mathbb{R}$ be defined as
\begin{equation}
q\left( k\right) :=\E_{\left\vert X\right\vert =k}\left[ p\left( X\right) 
\right] .
\end{equation}
Then $q$ can be written as a real polynomial in $k$ of degree at most $d$.
\end{lemma}

We now introduce a different, lesser known
notion of symmetrization, 
which we call the \emph{erase-all-subscripts} symmetrization
for reasons to be explained shortly. This symmetrization previously appeared in \cite{shi} under the name ``linearization,'' and it is also equivalent to the noise operator used in analysis of Boolean functions \cite[Definition 2.46]{odonnell}.

\begin{lemma}[Erase-all-subscripts
symmetrization]
Let $p:\left\{ 0,1\right\} ^{N}\rightarrow\mathbb{R}$\ be a
real multilinear polynomial of degree $d$,
and for
any real number $k \in [0, 1]$,
let $M_{k}$ denote
the distribution over $\{0, 1\}^N$, wherein each coordinate is selected
independently to be 1 with probability $k$.
Let $q:[0, 1] \to
\mathbb{R}$ be defined as
\begin{equation}
q\left( k\right) :=\E_{X \sim M_{k}}\left[ p\left( X\right) 
\right] .
\end{equation}
Then $q$ can be written as a real polynomial in $k$ of degree at most $d$.
\label{lem:eraseallsubscripts}
\end{lemma}
\begin{proof} (see, for example, \cite[Proof of Theorem 3]{STT12}).
Given the multivariate polynomial expansion of $p$,
we can obtain $q$ easily just by ``erasing all the subscripts in each variable''. For example, if
$p(x_1, x_2, x_3)= 2 x_1 x_2 + x_2 x_3 + x_2$, we replace every $x_i$ with $k$ to obtain
$q(k)=2k\cdot k+ k\cdot k + k=3k^2 + k$.
This follows from linearity of expectation along with the fact that $M_k$
is defined to be the
product distribution wherein
each coordinate has expected value $k$.
\end{proof}

We highlight the following key difference between Minsky--Papert symmetrization
and the erase-all-subscripts symmetrization. Let $p:\left\{ 0,1\right\} ^{N}\rightarrow [0, 1]$
be a real multivariate polynomial whose evaluations at Boolean inputs are in
$[0, 1]$, i.e., for all $x\in\{0,1\}^n$, we have $p(x)\in[0,1]$.
If $q$ is the erase-all-subscripts symmetrization of $p$, 
then $q$ takes values in $[0,1]$ at all \emph{real-valued}
inputs in $[0, 1]$: $q(k) \in [0,1]$ for all $k \in [0,1]$. 
If $q$ is the Minsky--Papert symmetrization of $p$,
then it is only guaranteed to take values in $[0,1]$ at 
\emph{integer-valued} inputs in $[0, N]$, i.e., $q(k)\in [0,1]$ 
is only guaranteed to hold at $k\in\{0,1,\ldots,N\}$.
This is the main reason we use erase-all-subscripts symmetrization
in this work.

\para{Bivariate symmetrizations.}
In this paper, it will be convenient
to consider bivariate versions of both
Minsky--Papert and erase-all-subscripts
symmetrization, and their applications to oracle separations. To this end,
define $X\in \left\{ 0,1\right\} ^{N}$, the \textquotedblleft characteristic
string\textquotedblright\ of the set $S\subseteq \left[ N\right] $, by $
x_{i}=1$\ if $i\in S$\ and $x_{i}=0$\ otherwise. \ Let $\mathcal{O}_{S}$
denote the unitary that performs a membership query to $S$, defined as
\begin{equation}
\mathcal{O}_{S}\left\vert i\right\rangle \left\vert b\right\rangle
=(1-2bx_{i})\left\vert i\right\rangle \left\vert b\right\rangle
\end{equation}
for any index $i\in \lbrack N]$ and bit $b\in \{0,1\}$.

Because we study oracle intersection problems, it is often convenient to
think of an algorithm as having access to \textit{two} oracles, wherein the
first bit in the oracle register selects the choice of oracle. \ As a
consequence, we need a slight generalization of a now well-established fact
in quantum complexity: that the acceptance probability of a quantum
algorithm with an oracle can be expressed as a polynomial in the bits of the
oracle string.

\begin{lemma}[Symmetrization with two oracles]
\label{lem:symmetrization} Let $Q^{\mathcal{O}_{S_{0}},\mathcal{O}_{S_{1}}}$
be a quantum algorithm that makes $T$ queries to a pair of membership
oracles for sets $S_{0},S_{1}\subseteq \lbrack N]$. \ Let $D_{\mu }$ denote
the distribution over subsets of $[N]$ wherein each element is selected
independently with probability $\frac{\mu }{N}$. \ Then there exist
bivariate real polynomials $q(s,t)$ and $p(x,y)$ of degree at most $2T$
satisfying:
\begin{align*}
\textrm{for all real numbers } s, t \in [0, N],& \quad
q(s,t)=\E_{\substack{ S_{0}\sim D_{s}, \\ S_{1}\sim D_{t}}}\left[ \Pr [Q^{
\mathcal{O}_{S_{0}},\mathcal{O}_{S_{1}}}\text{ accepts}]\right], \textrm{ and} \\
\textrm{for all integers } x, y \in \{0, 1, \dots, N\},& \quad
p(x,y)=\E_{\substack{ |S_{0}|=x, \\ |S_{1}|=y}}\left[ \Pr [Q^{\mathcal{O}
_{S_{0}},\mathcal{O}_{S_{1}}}\text{ accepts}]\right].
\end{align*}
\end{lemma}

\begin{proof}
Take $X = X_0|X_1$ to be the concatenation of the characteristic strings of the two oracles, and let $S \subseteq [2N]$ be such that $X$ is the characteristic string of $S$. Then, Lemma 4.2 of Beals et al.\ \cite{bbcmw} tells us that there is a real multilinear polynomial $r(X)$ of degree at most $2T$ in the bits of $X$ such that $r(X) = \Pr[ Q^{\mathcal{O}_S} \text{ accepts}]$.

Observe that $r$ has a meaningful probabilistic interpretation over arbitrary inputs in $[0, 1]$. A vector $X \in [0, 1]^{2N}$ of probabilities corresponds to a distribution over $\{0,1\}^{2N}$ wherein each bit is chosen from a Bernoulli distribution with the corresponding probability. Because $r$ is multilinear, $r$ in fact computes the expectation of the acceptance probability over this distribution. In particular, the polynomial
\begin{equation}
q(s, t) = r\bigg(\underbrace{\frac{s}{N},\ldots,\frac{s}{N}}_{N \text{ times}}, \underbrace{\frac{t}{N},\ldots,\frac{t}{N}}_{N \text{ times}}\bigg) = \E_{\substack{S_0 \sim D_s,\\S_1 \sim D_t}}\left[\Pr[ Q^{\mathcal{O}_{S_0},\mathcal{O}_{S_1}} \text{ accepts}] \right]    
\end{equation}
corresponds to selecting $S_0 \sim D_s$ and $S_1 \sim D_t$. The total degree of $q$ is obviously
at most the degree of $r$, by the same reasoning as in the proof of \Cref{lem:eraseallsubscripts}.

To construct $p$, we apply the symmetrization lemma of Minsky and Papert \cite{mp} to symmetrize $r$, first with respect to $X_0$, then with respect to $X_1$:
\begin{equation}
    p_0(x, X_1) = \E_{|S_0|=x} r(X_0,X_1) =  \E_{|S_0|=x}\left[\Pr[ Q^{\mathcal{O}_{S_0},\mathcal{O}_{S_1}} \text{ accepts}] \right]
\end{equation}
\begin{equation}
    p(x, y) = \E_{|S_1|=y} p_0(x,X_1) = \E_{\substack{|S_0|=x,\\|S_1|=y}}\left[\Pr[ Q^{\mathcal{O}_{S_0},\mathcal{O}_{S_1}} \text{ accepts}] \right] 
\end{equation}  
The degree of $p$ is at most the degree of $r$, due to \Cref{symlem}.
\end{proof}


We remark that, as a consequence of their definitions in \Cref{lem:symmetrization}, $p$ and $q$ satisfy:
\begin{equation}
q(s, t) = \E \left[ p(X, Y) \right],
\end{equation}
where $X$ and $Y$ are drawn from $N$-trial binomial distributions with means $s$ and $t$, respectively.

\para{Symmetric Laurent polynomials.}
Finally, we state a useful fact about Laurent polynomials:

\begin{lemma}[Symmetric Laurent polynomials]
\label{lem:symmetric} Let $\ell (x)$ be a real Laurent polynomial of positive
and negative degree $d$ that satisfies $\ell (x)=\ell (1/x)$. \ Then there exists a (ordinary) real
polynomial $q$ of degree $d$ such that $\ell (x)=q(x+1/x)$.
\end{lemma}

\begin{proof}
$\ell(x) = \ell(1/x)$ implies that the coefficients of the $x^i$ and $x^{-i}$ terms are equal for all $i$, as otherwise $\ell(x) - \ell(1/x)$ would not equal the zero polynomial. Thus, we may write $\ell(x) = \sum_{i=0}^d a_i \cdot (x^i + x^{-i})$ for some coefficients $a_i$. So, it suffices to show that $x^i + x^{-i}$ can be expressed as a polynomial in $x + 1/x$ for all $0 \le i \le d$.

We prove by induction on $i$. The case $i = 0$ corresponds to constant polynomials. For $i > 0$, by the binomial theorem, observe that $(x + 1/x)^i = x^i + x^{-i} + r(x)$ where $r$ is a degree $i - 1$ real Laurent polynomial satisfying $r(x) = r(1/x)$. By the induction assumption, $r$ can be expressed as a polynomial in $x + 1/x$, so we have $x^i + x^{-i} = (x + 1/x)^i - r(x)$ is expressed as a polynomial in $x + 1/x$.
\end{proof}

%
\subsection{Complexity classes}

\begin{definition}
\label{def:qma}
The complexity class $\mathsf{QMA}$ consists of the
languages $L$ for which there exists a quantum polynomial time
verifier $V$ with the following properties:

\begin{enumerate}
\item Completeness: if $x \in L$, then there exists a quantum witness state $|\psi\>$ on $\mathrm{poly}(|x|)$ qubits such that $\Pr\left[V(x, |\psi\>) \text{ accepts}\right] \ge \frac{2}{3}$.

\item Soundness: if $x \not\in L$, then for any quantum witness state $|\psi\>$ on $\mathrm{poly}(|x|)$ qubits, $\Pr\left[V(x, |\psi\>) \text{ accepts}\right] \le \frac{1}{3}$.
\end{enumerate}
\end{definition}

A quantum verifier that satisfies the above
promise for a particular language will be referred to as a $\mathsf{QMA}$
verifier or $\mathsf{QMA}$ protocol throughout.

Though $\mathsf{SBP}$ and $\mathsf{SBQP}$ can be defined in terms of
counting complexity functions, for our purposes it is easier to
work with the following equivalent definitions (see B\"ohler et al.\ \cite
{BGM06}):

\begin{definition}
\label{def:sbp_sbqp} The complexity class $\mathsf{SBP}$ consists of the
languages $L$ for which there exists a probabilistic polynomial time
algorithm $M$ and a polynomial $\sigma$ with the following properties:

\begin{enumerate}
\item If $x \in L$, then $\Pr\left[M(x) \text{ accepts}\right] \ge
2^{-\sigma(|x|)}$.

\item If $x \not\in L$, then $\Pr\left[M(x) \text{ accepts}\right] \le
2^{-\sigma(|x|)}/2$.
\end{enumerate}

The complexity class $\mathsf{SBQP}$ is defined analogously, wherein the
classical algorithm is replaced with a quantum algorithm.
\end{definition}

A classical (respectively, quantum) algorithm that satisfies the above
promise for a particular language will be referred to as an $\mathsf{SBP}$
(respectively, $\mathsf{SBQP}$) algorithm throughout. \ Using these
definitions, a  query complexity relation between $\mathsf{QMA}$
protocols and $\mathsf{SBQP}$ algorithms follows from the procedure of
Marriott and Watrous \cite{marriott}, which shows that one can exponentially
improve the soundness and completeness errors of a $\mathsf{QMA}$ protocol
without increasing the witness size. \ 
This relationship is now standard; see
for example \cite[Remark 6]{marriott} or \cite[Proposition 4.2]
{sherstov2019vanishing} for a proof of the following lemma:

\begin{lemma} 
\label{lem:guessing} Suppose there is a $\mathsf{QMA}$ protocol for some
problem that makes $T$ queries and receives an $m$-qubit witness. \ Then
there is a quantum query algorithm $Q$ for the same problem that makes $O(mT)$ queries, and satisfies the following:

\begin{enumerate}
\item If $x \in L$, then $\Pr\left[Q(x) \text{ accepts}\right] \ge
2^{-m}$.

\item If $x \not\in L$, then $\Pr\left[Q(x) \text{ accepts}\right] \le
2^{-10m}$.
\end{enumerate}
\end{lemma}


%
